\documentclass[11pt]{article}
\pagestyle{plain}

% \pgfplotsset{compat=1.18}
\usepackage[margin=1in]{geometry}
\usepackage{amsmath,amsfonts,amssymb,amsthm}
\usepackage[ruled,vlined,linesnumbered,commentsnumbered]{algorithm2e}
\SetKwInput{KwIn}{Input}
\SetKwInput{KwOut}{Output}
\SetKwInput{KwGlobal}{Global variable}
\SetKwInput{KwOracle}{Oracle}
\SetKwComment{Comment}{// }{}
\usepackage{graphicx}
\usepackage{enumerate}
\usepackage{bbm}
\usepackage{verbatim}
\usepackage{hyperref,color}
\usepackage[capitalize,nameinlink]{cleveref}
\usepackage[dvipsnames]{xcolor}
\hypersetup{
	colorlinks=true,
	pdfpagemode=UseNone,
	citecolor=OliveGreen,
	linkcolor=NavyBlue,
	urlcolor=Magenta,
	pdfstartview=FitW
}
\usepackage{appendix}
\usepackage{makecell}
\usepackage{tablefootnote}
\crefname{appsec}{Appendix}{Appendices}
\usepackage{tikz}
\usepackage{pgfplots}
\usepackage{xifthen}
\usepackage{subcaption}
\usepackage{hhline}

\theoremstyle{plain}
\newtheorem{theorem}{Theorem}[section]
\newtheorem{proposition}[theorem]{Proposition}
\newtheorem{lemma}[theorem]{Lemma}
\newtheorem{corollary}[theorem]{Corollary}
\newtheorem{conjecture}[theorem]{Conjecture}
\newtheorem{problem}[theorem]{Problem}
\newtheorem{claim}[theorem]{Claim}
\newtheorem{fact}[theorem]{Fact}

\theoremstyle{definition}
\newtheorem{definition}[theorem]{Definition}
\newtheorem{example}[theorem]{Example}
\newtheorem{observation}[theorem]{Observation}
\newtheorem{assumption}[theorem]{Assumption}
\newtheorem*{assumption*}{Assumption}
\newtheorem{condition}[theorem]{Condition}

\theoremstyle{remark}
\newtheorem{remark}[theorem]{Remark}

\crefname{lemma}{Lemma}{Lemmas}
\crefname{theorem}{Theorem}{Theorems}
\crefname{definition}{Definition}{Definitions}
\crefname{fact}{Fact}{Facts}
\crefname{claim}{Claim}{Claims}
\crefname{proposition}{Proposition}{Propositions}


\newcommand{\dif}{\,\mathrm{d}}
%\newcommand{\E}{\mathbb{E}}
%\newcommand{\Var}{\mathrm{Var}}
\newcommand{\Cov}{\mathrm{Cov}}
\newcommand{\MI}{\mathrm{I}}
%\newcommand{\Ent}{\mathrm{Ent}}
\newcommand{\ind}{\mathbbm{1}}
\newcommand{\allone}{\mathbf{1}}
\newcommand{\tr}{\mathrm{tr}}
\newcommand{\wt}{\mathrm{weight}}
\DeclareMathOperator*{\argmax}{arg\,max}
\newcommand{\norm}[1]{\left\lVert #1 \right\rVert}
\newcommand{\TV}[2]{\left\|#1 - #2\right\|_{\mathrm{TV}}}
\newcommand{\ceil}[1]{\left\lceil #1 \right\rceil}
\newcommand{\floor}[1]{\left\lfloor #1 \right\rfloor}
\newcommand{\ip}[2]{\left\langle #1 , #2 \right\rangle}
\newcommand{\grad}{\nabla}
\newcommand{\hessian}{\nabla^2}
\newcommand{\trans}{\intercal}
\newcommand{\diag}{\mathrm{diag}}
\newcommand{\poly}{\mathrm{poly}}
\newcommand{\dist}{\mathrm{dist}}
\newcommand{\sgn}{\mathrm{sgn}}
\newcommand{\fpras}{\mathsf{FPRAS}}
\newcommand{\fpaus}{\mathsf{FPAUS}}
\newcommand{\fptas}{\mathsf{FPTAS}}

\newcommand{\pll}{\parallel}
\newcommand{\eps}{\varepsilon}

\newcommand{\N}{\mathbb{N}}
\newcommand{\Z}{\mathbb{Z}}
\newcommand{\Q}{\mathbb{Q}}
\newcommand{\R}{\mathbb{R}}
\newcommand{\C}{\mathbb{C}}

\renewcommand{\AA}{\mathcal{A}}
\newcommand{\BB}{\mathcal{B}}
\newcommand{\CC}{\mathcal{C}}
\newcommand{\DD}{\mathcal{D}}
\newcommand{\EE}{\mathcal{E}}
\newcommand{\FF}{\mathcal{F}}
\newcommand{\GG}{\mathcal{G}}
\newcommand{\HH}{\mathcal{H}}
\newcommand{\II}{\mathcal{I}}
\newcommand{\JJ}{\mathcal{J}}
\newcommand{\KK}{\mathcal{K}}
\newcommand{\LL}{\mathcal{L}}
\newcommand{\MM}{\mathcal{M}}
\newcommand{\NN}{\mathcal{N}}
\newcommand{\OO}{\mathcal{O}}
\newcommand{\PP}{\mathcal{P}}
\newcommand{\QQ}{\mathcal{Q}}
\newcommand{\RR}{\mathcal{R}}
\renewcommand{\SS}{\mathcal{S}}
\newcommand{\TT}{\mathcal{T}}
\newcommand{\UU}{\mathcal{U}}
\newcommand{\VV}{\mathcal{V}}
\newcommand{\WW}{\mathcal{W}}
\newcommand{\XX}{\mathcal{X}}
\newcommand{\YY}{\mathcal{Y}}
\newcommand{\ZZ}{\mathcal{Z}}

\newcommand{\KL}[2]{D_{\mathrm{KL}} \left( #1 \,\Vert\, #2 \right)}
\newcommand{\KLbig}[2]{D_{\mathrm{KL}} \left( #1 \,\big\Vert\, #2 \right)}
\newcommand{\KLBig}[2]{D_{\mathrm{KL}} \left( #1 \,\Big\Vert\, #2 \right)}

\newcommand{\SKL}[2]{D_{\mathrm{SKL}} \left( #1 \,\Vert\, #2 \right)}
\newcommand{\SKLbig}[2]{D_{\mathrm{SKL}} \left( #1 \,\big\Vert\, #2 \right)}
\newcommand{\SKLBig}[2]{D_{\mathrm{SKL}} \left( #1 \,\Big\Vert\, #2 \right)}

\newcommand{\qandq}{\quad\text{and}\quad}
\newcommand{\supp}{\mathrm{supp}}

\newcommand{\DKL}[2]{\-D_{\-{KL}}\left(#1 \parallel #2\right)}
\newcommand{\DTV}[2]{\-D_{\mathrm{TV}}\left({#1},{#2}\right)}
% \newcommand{\dist}{\mathrm{dist}}
\newcommand{\e}{\mathrm{e}}
\renewcommand{\epsilon}{\varepsilon}
\renewcommand{\emptyset}{\varnothing}
% \newcommand{\norm}[1]{\left\Vert#1\right\Vert}
\newcommand{\set}[1]{\left\{#1\right\}}
\newcommand{\tuple}[1]{\left(#1\right)} 
% \newcommand{\eps}{\varepsilon}
\newcommand{\inner}[2]{\left\langle #1,#2\right\rangle}
\newcommand{\tp}{\tuple}
\newcommand{\ol}{\overline}
\newcommand{\abs}[1]{\left\vert#1\right\vert}
\newcommand{\ctp}[1]{\left\lceil#1\right\rceil}
\newcommand{\ftp}[1]{\left\lfloor#1\right\rfloor}
\newcommand{\Ind}[1]{\-{Ind}_{#1}}
\newcommand{\todo}[1]{{\color{red} TODO: #1}}

\newcommand{\Cor}{\Psi^{\mathrm{sym}}}
\newcommand{\cor}{\Psi^{\-{cor}}}
\newcommand{\Inf}{\Psi}
% \newcommand{\diag}{\-{diag}}

\def\*#1{\boldsymbol{#1}} % Use \*A for \mathbf{A}
\def\+#1{\mathcal{#1}} % Use \+A for \mathcal{A}
\def\-#1{\mathrm{#1}} % Use \-A for \mathrm{A}
\def\=#1{\mathbb{#1}} % Use \^A for \mathbb{A}
\def\!#1{\mathfrak{#1}} % Use \!A for \mathfrak{A}

\newcommand*\diff{\mathop{}\!\mathrm{d}}

% probability
% \DeclareMathOperator*{\oPr}{\mathbf{Pr}}
\def\oPr{\mathop{\mathrm{Pr}}}
\renewcommand{\Pr}[2][]{ \ifthenelse{\isempty{#1}}
  {\oPr\left[#2\right]}
  {\oPr_{#1}\left[#2\right]} } % Use \Pr[a]{b} for \mathbf{Pr}_a[b], \Pr{b} for  \mathbf{Pr}[b]

% expectation: relies on the package "dsfont"
% \DeclareMathOperator*{\oE}{\mathbb{E}}
\def\oE{\mathop{\mathbb{E}}}
\newcommand{\E}[2][]{ \ifthenelse{\isempty{#1}}
  {\oE\left[#2\right]}
  {\oE_{#1}\left[#2\right]} }

% variance
%\DeclareMathOperator*{\oVar}{\mathbf{Var}}
\def\oVar{\mathrm{Var}}
\newcommand{\Var}[2][]{ \ifthenelse{\isempty{#1}}
  {\oVar\left[#2\right]}
  {\oVar_{#1}\left[#2\right]} }

% entropy
% \DeclareMathOperator*{\oEnt}{\mathbf{Ent}}
\def\oEnt{\mathrm{Ent}}
\newcommand{\Ent}[2][]{ \ifthenelse{\isempty{#1}}
  {\oEnt\left[#2\right]}
  {\oEnt_{#1}\left[#2\right]} }

\newcommand{\PhiEnt}[2][]{ \ifthenelse{\isempty{#1}}
  {\oEnt^\Phi\left[#2\right]}
  {\oEnt^\Phi_{#1}\left[#2\right]} }


\newenvironment{Exampleparagraph}[1][]
{
    \par 
    \noindent 
    \textbf{\color{blue} #1}
    \par 
    \vspace{0.5em} 
    \itshape 
}
{
    \par 
    \vspace{1em} 
}

%\input{local-macro}

\newcommand{\RED}[1]{\textcolor{red}{#1}}
\newcommand{\zc}[1]{\textcolor{red}{ZC: #1}}

\newcommand{\bern}[1]{\mathrm{Bern}\left(#1\right)}




\title{Second-order $\Phi$-Entropic Stability}
\author{}
\date{}

\begin{document}

\maketitle

Recall that $\mu_1 = \mu$ and $\mu_0 = \delta_\phi$.
We define
\begin{align*}
    G(\theta) = \E[\tau \sim \mu_\theta]{ \PhiEnt[\mu^{\theta,\tau}_1]{f} }.
\end{align*}
Notice that
$G(1) = 0$ and $G(0) = \PhiEnt[\mu]{f} $.


\section{First-order $\Phi$-Entropic Stability}

\begin{definition}
For any $\theta \in [0, 1)$ and any pinning $\tau$, it holds
    \begin{align*}
        \lim_{h \to 0^+} \frac{1}{h} \PhiEnt[\sigma \sim \mu_{\theta+h}^{\theta,\tau}]{ \E[\mu_1^{\theta+h,\sigma}]{f} } 
        \le \frac{C(\theta)}{1-\theta} \PhiEnt[\mu_1^{\theta,\tau}]{f}.
    \end{align*}
\end{definition}

Observe that
\begin{align*}
    -G'(\theta) 
    &= \lim_{h \to 0^+} \frac{1}{h} \left( \E[\tau \sim \mu_{\theta}]{ \PhiEnt[\mu_1^{\theta,\tau}]{f}} - \E[\sigma \sim \mu_{\theta+h}]{ \PhiEnt[\mu_1^{\theta+h,\sigma}]{f} } \right) \\
    & = \lim_{h \to 0^+} \frac{1}{h} \E[\tau \sim \mu_\theta] { \PhiEnt[\sigma \sim \mu_{\theta+h}^{\theta,\tau}]{ \E[\mu_1^{\theta+h,\sigma}]{f} } } \\
    & = \E[\tau \sim \mu_\theta] { \lim_{h \to 0^+} \frac{1}{h} \PhiEnt[\sigma \sim \mu_{\theta+h}^{\theta,\tau}] {\E[\mu_1^{\theta+h,\sigma}]{f}} }.
\end{align*}

Therefore,
\begin{align*}
    & \text{First-order $\Phi$-Entropic Stability} \\
    \implies &
    -G'(\theta) \le \frac{C(\theta)}{1-\theta} G(\theta) \\
    \implies &
    G(0) \le \exp \left( \int^{\theta}_{0} \frac{C(\eta)}{1-\eta} \dif \eta \right) G(\theta)
\end{align*}


\section{Second-order $\Phi$-Entropic Stability}

\begin{definition}
For any $\theta \in [0, 1)$ and any pinning $\tau$, it holds
    \begin{multline*}
        \lim_{h \to 0^+} \frac{1}{h^2} \left( \PhiEnt[\mu_1^{\theta,\tau}]{f} - 2 \E[\sigma \sim \mu_{\theta+h}^{\theta,\tau}]{ \PhiEnt[\mu_1^{\theta+h,\sigma}]{f}} + \E[\rho \sim \mu_{\theta+2h}^{\theta,\tau}]{ \PhiEnt[\mu_1^{\theta+2h,\rho}]{f} } \right) \\
        \le \frac{C(\theta)}{1-\theta} \lim_{h \to 0^+} \frac{1}{h} \PhiEnt[\sigma \sim \mu_{\theta+h}^{\theta,\tau}]{ \E[\mu_1^{\theta+h,\sigma}]{f}}.
    \end{multline*}
\end{definition}

Observe that
\begin{align*}
    G''(\theta) = \E[\tau \sim \mu_\theta]{\lim_{h \to 0^+} \frac{1}{h^2} \left( \PhiEnt[\mu_1^{\theta,\tau}]{f} - 2 \E[\sigma \sim \mu_{\theta+h}^{\theta,\tau}]{ \PhiEnt[\mu_1^{\theta+h,\sigma}]{f}} + \E[\rho \sim \mu_{\theta+2h}^{\theta,\tau}] { \PhiEnt[\mu_1^{\theta+2h,\rho}]{f} } \right) }.
\end{align*}

Therefore,
\begin{align*}
    & \text{Second-order $\Phi$-Entropic Stability} \\
    \implies &
    G''(\theta) \le -\frac{C(\theta)}{1-\theta} G'(\theta) \\
    \implies &
    G(0) \le [\cdots]\; G(\theta)
\end{align*}

\subsection{Spectral Stability}
Define $f_\eta(x) = \E[Q_{\eta\to 1}(x,\cdot)]{f} = \E[\mu_{1}^{\eta,x}]{f}$ and $\mu_\eta$ to be the distribution of $Y_\eta$.
% \begin{definition}[Second order of $\phi$-Entropic Stability]
%     For any $\theta \in [0,1)$ and any pinning $\tau$, it holds
%     \begin{align*}
%         &\lim_{h\to 0^+} \frac 1{h^2}\left(\PhiEnt[\mu_1^{\theta,\tau}]{f} - 2\E[\sigma\sim \mu_{\theta+h}^{\theta,\tau}]{\PhiEnt[\mu_{1}^{\theta+h,\sigma}]{f}} + \E[\rho\sim \mu_{\theta+2h}^{\theta,\tau}]{\PhiEnt[\mu_{1}^{\theta+2h,\rho}]{f}}\right) \\\leq&C(\theta) \lim_{h\to 0^+}\frac 1h \PhiEnt[\sigma\sim \mu_{\theta+h}^{\theta,\tau}]{f_{\theta+h}(\sigma)} .
%     \end{align*}
% \end{definition}
We take $\phi = x^2$ and $P/Q$ be the down/up walk in field dynamics as example.
Thus, we have that 
$$
    \forall y\subseteq x, P_{1\to \theta}(x,y) = \theta^{y}(1-\theta)^{x-y}
$$
$$
    \forall y\subseteq x, Q_{\theta \to 1}(y,x) = \Pr{Y_1 = x\mid Y_\theta = y} = \theta^y(1-\theta)^{x-y} \frac{\Pr{Y_1=x}}{\Pr{Y_\theta=y}} \propto ((1-\theta)*\mu) (y).
$$
$$
    P_{\eta + h \to \eta}(x,y) = \left(\frac{\eta}{\eta+h}\right)^{y}\left(\frac{h}{\eta+h}\right)^{x-y} = \begin{cases}
        1-\frac{h|x|}{\eta} + o(h) &, x=y
        \\ \frac{h}{\eta} + o(h) &, y=x-1
        \\o(h) &, o.w.
    \end{cases}.
$$
$$
    Q_{\eta\to \eta+h}(x,y) = \Pr{Y_\eta = x\mid Y_{\eta+h} = y}\frac{\Pr{Y_{\eta+h} = y}}{\Pr{Y_{\eta} = x}} = \begin{cases}
        1 - \frac{h}{\eta}\sum_{i\notin x}\frac{\Pr{Y_\eta = x+i}}{\Pr{Y_\eta = x}} + o(h)&, x=y
        \\ \frac{\Pr{Y_\eta = y}}{\Pr{Y_\eta = x}}\cdot\frac{h}{\eta} + o(h) &, y - x = 1
        \\ o(h) &, o.w.
    \end{cases} .
$$
And we have that
\begin{align*}
    \Pr{Y_{\eta+h} = y} &= \sum_{x\subseteq y} \Pr{Y_{\eta} = x}Q_{\eta\to \eta+h}(x,y) 
    \\&= \Pr{Y_\eta = y} - \frac h\eta\sum_{i\notin y}\Pr{Y_\eta = y+i}+\frac h\eta\sum_{i\in y} \Pr{Y_{\eta} = y}  + o(h).
    \\&=\Pr{Y_\eta = y} - \frac h\eta\sum_{i\notin y}\Pr{Y_\eta = y+i}+\frac {h|y|}\eta\Pr{Y_{\eta} = y}  + o(h).
\end{align*}
Therefore we have the following lemma.
\begin{lemma}\label{lem:first_order}
We have that for $\eta \in [0,1)$ and pinning $x$,
$$
    \lim_{h\to 0^+} \frac 1h \Var[Q_{\eta\to \eta+h}(x,\cdot)]{f_{\eta+ h}} = \frac 1\eta \sum_{i\notin x} \frac{\mu_\eta(x+i)}{\mu_\eta(x)}(f_\eta(x+i) - f_\eta(x))^2 .
$$
\end{lemma}
\begin{proof}
    The following should be true,
    \begin{align*}
        \Var[Q_{\eta\to \eta+h}(x,\cdot)]{f_{\eta+h}} &= \sum_{x\subseteq y} Q_{\eta\to \eta+h}(x,y) (f_{\eta+h}(y) - \E[Q_{\eta\to \eta+h}(x,\cdot)]{f_{\eta+h}})^2
        \\&= o(h) + \sum_{i\notin x} Q_{\eta\to \eta+h}(x,x+i) (f_{\eta+h}(x+i) - f_{\eta}(x))^2
        \\&= o(h) + \frac{h}{\eta} \sum_{i\notin x} \frac{\Pr{Y_\eta = x+i}}{\Pr{Y_\eta = x}}(f_{\eta}(x+i) - f_{\eta}(x))^2, 
    \end{align*}
    where $f_{\eta+h}(x) = (1+O(h))f_\eta(x)$. 
\end{proof}
For fixed pinning $x$, we want to compute
\begin{align*}
    \frac{1}{(\eta + h)\eta}\left(\eta\sum_{i\notin x} \frac{\Pr{Y_{\eta+h} = x+i}}{\Pr{Y_{\eta+h} = x}}(f_{\eta+h}(x+i) - f_{\eta+h}(x))^2 - (\eta+h)\sum_{i\notin x} \frac{\Pr{Y_\eta = x+i}}{\Pr{Y_\eta = x}}(f_\eta(x+i) - f_\eta(x))^2\right).
\end{align*}
% $$
%     \Pr{Y_{\eta+h} = x} = \sum_{y\subseteq x} \Pr{Y_{\eta} = y}Q_{\eta\to \eta+h}(y,x)
% $$
% The first term is that
% $$
%     -\frac{h}{\eta^2}\sum_{i\notin x} \frac{\Pr{Y_{\eta} = x+i}}{\Pr{Y_{\eta} = x}}(f_{\eta}(x+i) - f_{\eta}(x))^2 .
% $$
% The second term is that
% $$
%     \frac{h}{\eta}
% $$
The second derivative is that
\begin{align*}
    \frac 1\eta\sum_{i\notin x}\left([h]\frac{\Pr{Y_{\eta+h} = x+i}}{\Pr{Y_{\eta+h} = x}}\right)(f_{\eta}(x+i) - f_{\eta}(x))^2 &+ \sum_{i\notin x}\frac 1\eta\left(\frac{\Pr{Y_{\eta} = x+i}}{\Pr{Y_{\eta} = x}}\right)[h](f_{\eta+h}(x+i) - f_{\eta+h}(x))^2
    \\& - \frac 1{\eta^2}\sum_{i\notin x} \frac{\Pr{Y_\eta = x+i}}{\Pr{Y_\eta = x}}(f_\eta(x+i) - f_\eta(x))^2 .
\end{align*}
Note that 
\begin{align*}
    Q_{\eta+h \to 1}(x,y) &= \sum_{x\subseteq z} Q_{\eta\to \eta+h}^{-1}(x,z)Q_{\eta\to 1}(z,y)
    \\&= Q_{\eta\to \eta+h}^{-1}(x,x)Q_{\eta\to 1}(x,y) + \sum_{i\notin x} Q_{\eta\to \eta+h}^{-1}(x,x+i)Q_{\eta\to 1}(x+i,y)  + o(h)
    \\&= Q_{\eta\to 1}(x,y) + \frac h\eta\sum_{i\notin x} \frac{\Pr{Y_{\eta} = x+i}}{\Pr{Y_\eta = x}} (Q_{\eta\to 1}(x,y) - Q_{\eta\to 1}(x+i,y)) + o(h),
\end{align*}
since $(I+Lh+o(h))^{-1} = I - Lh + o(h)$.
Therefore, 
\begin{align*}
    f_{\eta+h}(x) - f_\eta(x) &= \sum_{x\subseteq y} (Q_{\eta+h\to 1} - Q_{\eta\to 1} )(x,y)f(y) 
    \\&= \frac{h}{\eta}\sum_{y:x\subseteq y}\sum_{i\notin x} \frac{\Pr{Y_{\eta} = x+i}}{\Pr{Y_\eta = x}} (Q_{\eta\to 1}(x,y) - Q_{\eta\to 1}(x+i,y)) f(y)  + o(h)
    \\&= \frac{h}{\eta}\sum_{i\notin x}\frac{\Pr{Y_{\eta} = x+i}}{\Pr{Y_\eta = x}}(\sum_{y:x\subseteq y}Q_{\eta\to 1}(x,y) f(y) - \sum_{y:x\subseteq y}  Q_{\eta\to 1}(x+i,y) f(y) )  + o(h),
    \\&= \frac{h}{\eta}\sum_{i\notin x}\frac{\Pr{Y_{\eta} = x+i}}{\Pr{Y_\eta = x}}(f_\eta(x) - f_\eta(x+i)) + o(h),
\end{align*}
Thus, 
\begin{align*}
    [h](f_{\eta+h}(x+i) - f_{\eta+h}(x))^2 &= \frac {2}\eta(f_{\eta}(x+i) - f_{\eta}(x))(\sum_{j\notin x+i}\frac{\Pr{Y_{\eta} = x+j+i}}{\Pr{Y_\eta = x+i}}(f_\eta(x+i) - f_\eta(x+j+i)) 
    \\&\qquad\qquad\qquad\qquad\qquad - \sum_{j\notin x}\frac{\Pr{Y_{\eta} = x+j}}{\Pr{Y_\eta = x}}(f_\eta(x) - f_\eta(x+j))) .
\end{align*}
For the first term, since $\frac{A+Bh}{C+Dh} =\frac{A}{C} +  \frac{CB-AD}{C^2}h+o(h)$,
\begin{align*}
    [h]\frac{\Pr{Y_{\eta+h} = x+i}}{\Pr{Y_{\eta+h} = x}} &= -\frac{\Pr{Y_{\eta} = x }\Pr{Y_{\eta} = x+i}}{\eta\Pr{Y_\eta = x}^2}\Bigg(\sum_{j\notin x+i}\frac{\Pr{Y_{\eta} = x+j+i}}{\Pr{Y_\eta = x+i}} - |x+i|
    \\&\qquad\qquad\qquad\qquad\qquad\qquad\qquad-\left(\sum_{j\notin x}\frac{\Pr{Y_{\eta} = x+j}}{\Pr{Y_\eta = x}} - |x|\right)\Bigg)
    \\&= \frac{\Pr{Y_{\eta} = x+i}}{\eta\Pr{Y_\eta = x}}\left(1-\sum_{j\notin x+i} \frac{\Pr{Y_\eta= x+i+j}}{\Pr{Y_\eta = x+i}} + \sum_{j\notin x}\frac{\Pr{Y_\eta = x+j}}{\Pr{Y_\eta = x}}\right)&
\end{align*}
% $$
%     \Pr{Y_{\eta+h} = x + i} =  \Pr{Y_\eta = x + i} + \sum_{j\in x + i} \frac{h\Pr{Y_{\eta} = x + i}}{\eta\Pr{Y_{\eta} = x + i - j}}(\Pr{Y_\eta = x+i-j} - \Pr{Y_\eta = x+i}).
% $$
% $$
%     \Pr{Y_{\eta+h} = x } =  \Pr{Y_\eta = x } + \sum_{j\in x } \frac{h\Pr{Y_{\eta} = x }}{\eta\Pr{Y_{\eta} = x  - j}}(\Pr{Y_\eta = x-j} - \Pr{Y_\eta = x}).
% $$


In conclusion, the second derivative is that
\begin{align*}
    \frac{1}{\eta^2}\sum_{i\notin x} \frac{\Pr{Y_\eta = x+i}}{\Pr{Y_\eta = x}}(f_\eta(x+i) - f_\eta(x))^2(&\left(1-\sum_{j\notin x+i} \frac{\Pr{Y_\eta= x+i+j}}{\Pr{Y_\eta = x+i}} + \sum_{j\notin x}\frac{\Pr{Y_\eta = x+j}}{\Pr{Y_\eta = x}}\right)
    \\&+2\Bigg(\sum_{j\notin x+i}\frac{\Pr{Y_{\eta} = x+j+i}}{\Pr{Y_\eta = x+i}}\frac{f_\eta(x+i) - f_\eta(x+j+i)}{f_\eta(x+i) - f_\eta(x)} 
    \\&\qquad\ \  + \sum_{j\notin x}\frac{\Pr{Y_{\eta} = x+j}}{\Pr{Y_\eta = x}}\frac{f_\eta(x+j) - f_\eta(x)}{f_\eta(x+i) - f_\eta(x)} \Bigg)
    \\&- 1)
\end{align*}
That is, after simplification, the second derivative is that
\begin{align*}
\frac{1}{\eta^2}\sum_{i\notin x} \frac{\Pr{Y_\eta = x+i}}{\Pr{Y_\eta = x}}(f_\eta(x+i) - f_\eta(x))\Bigg(&\sum_{j\notin x+i}\frac{\Pr{Y_\eta= x+i+j}}{\Pr{Y_\eta = x+i}}\left(f_\eta(x)+f_\eta(x+i)-2f_\eta(x+i+j)\right)
\\&+\sum_{j\notin x}\frac{\Pr{Y_{\eta} = x+j}}{\Pr{Y_\eta = x}}\left(2f_\eta(x+j)+f_\eta(x+i)-3f_\eta(x)\right)\Bigg)
\end{align*}
Let $\Delta_i(x) = f_\eta(x+i) - f_\eta(x)$, $r_i(x) = \frac{\Pr{Y_\eta = x+i}}{\Pr{Y_\eta = x}}$, $S(x) = \sum_{i\notin x} r_i(x)$ and $T(x) = \sum_{i\notin x} r_i(x)\Delta_i(x)$, then the second derivative is that
\begin{equation}\label{eq:second_derivative}
    \frac 1{\eta^2}\sum_{i\notin x} r_i(x)\left(\Delta_i(x)^2(S(x)-S(x+i)) + 2\Delta_i(x)(T(x)-T(x+i))\right) .
\end{equation}
Note that the first derivative is that
$$
\frac 1\eta \sum_{i\notin x} \frac{\Pr{Y_\eta = x+i}}{\Pr{Y_\eta = x}}(f_\eta(x+i) - f_\eta(x))^2 = \frac 1\eta \sum_{i\notin x} r_i(x)\Delta_i(x)^2 .
$$
\begin{example}[$\eta = 0$]
    When $\eta = 0$, only $x = \emptyset$ makes sense.
    We omit $(\emptyset)$ for simplicity.
    We have that
    $$
        \lim_{\eta\to 0}\frac{r_i}\eta = \lim_{\eta \to 0}
        \frac {\eta\mu(\{i\})}{\eta Z(\eta*\mu)} = \frac{\mu(\{i\})}{\mu(\emptyset)},
    $$
    since $\Pr{Y_0 = \emptyset} = 1$.
    Thus, Equation \ref{eq:second_derivative} is as follows,
    \begin{align*}
        \sum_{i=1}^n \frac{\mu(\{i\})}{\mu(\emptyset)}\Bigg(&\Delta_i(\emptyset)^2\left(\frac{\mu(\{i\})}{\mu(\emptyset)} + \sum_{j\neq i}\frac{\mu(\{j\})^2 - \mu(\{i,j\})\mu(\emptyset)}{\mu(\emptyset)\mu(\{j\})}\right) 
        \\&+ 2\Delta_i(\emptyset)\left(\frac{\mu(\{i\})\Delta_i(\emptyset)}{\mu(\emptyset)} + \sum_{j\neq i}\frac{\mu(\{j\})^2\Delta_j(\emptyset) - \mu(\{i,j\})\mu(\emptyset)\Delta_j(\{i\})}{\mu(\emptyset)\mu(\{j\})}\right)\Bigg) 
    \end{align*}
\end{example}


\subsection{$\*{f}$ is a linear function}
Let $f(x) = \inner{\alpha}{x}$, we have that 
$$
    f_\eta(x) = \E[Q_{\eta\to 1}(x,y)]{\inner{\alpha}{y}} = \inner{\*m((1-\eta) * \mu_x)}{\alpha} .
$$
Let $\nu_x = Q_{\eta\to 1}(x,\cdot) = (1-\eta)*\mu_x$.
Since for $i\notin x$,
$$
    \nu_x(y\mid i\in y) = \frac{\ind [x+i\subseteq y]\nu_x(y)}{\*m(\nu_x)_i} \propto \ind [x+i\subseteq y](1-\eta)^{|y|}\mu(y) = \nu_{x+i}(y),
$$
Therefore, 
$$
    \Delta_i(x) = \sum_{j\notin x}(\nu_x(j\in y\mid i\in y) - \nu_x(j\in y))\alpha_j = \frac{1}{\nu_x(i\in y)} \inner{\Cov(\nu_x)_i}{\alpha} .
$$
By the law of $Y_\eta$, we have that 
$$
    r_i(x) = \frac{\Pr{Y_\eta=x+i}}{\Pr{Y_\eta = x}}  = \frac{\sum_{x+i\subseteq y} \mu(y) \eta^{|x+i|} (1-\eta)^{|y|-|x+i|}}{\sum_{x\subseteq y} \mu(y) \eta^{|x|} (1-\eta)^{|y|-|x|}} = \frac{\eta}{1-\eta} \nu_x(i\in y).
$$
Plugging all the above, the first derivative is as follows,
$$
    \frac 1{1 - \eta} \alpha^\top \Cov(\nu_x)\diag\{\*m(\nu_x)\}^{\dagger} \Cov(\nu_x)\alpha^\top 
$$
\begin{proposition}
    If $f$ is a pure linear function, then $\lim_{h \to 0^+} \frac{1}{h} \Var[Q_{\eta \to \eta + h}(x,\cdot)]{f_{\eta + h}} \le \frac{C(\eta)}{1-\eta} \cdot \Var[Q_{\eta \to 1}(x,\cdot)]{f}$ is equivalent to $\Cov(\nu_x)\preceq C(\eta)\cdot \diag\{\*m(\nu_x)\}$ .
\end{proposition}
\begin{proof}
    We complete the proof by the following equality.
    \begin{align*}
        \Var[Q_{\eta \to 1}(x,\cdot)]{f} &= \E[y\sim \nu_x]{\inner{\alpha}{y}^2} - \inner{\*m(\nu_x)}{\alpha}^2
       \\&= \alpha^\top \left(\E[y\sim \nu_x]{yy^\top} - \*m(\nu_x)\*m(\nu_x)^\top \right)\alpha = \alpha^\top \Cov(\nu_x)\alpha . \qedhere
    \end{align*}
\end{proof}
Following the same way, we have that
$$
    T(x) = \frac{\eta}{1-\eta}\*1^\top \Cov(\nu_x)\alpha, S(x) = \frac{\eta}{1-\eta}\E[y\sim \nu_x]{|y| - |x|} .
$$
Define 
$$
    M_x = \begin{pmatrix}
    \vdots\\\*1^\top (I-\Cov(\nu_{x+i})\Cov(\nu_x)^{-1})\\
    \vdots
    \end{pmatrix}_{i\notin x} .
$$
Then the second derivative is as follows,
\begin{align*}
    \alpha^\top \frac1{(1-\eta)^2}\Bigg(&\Cov(\nu_x)\diag\{\*m(\nu_x)\}^\dagger\diag\{\E[y\sim \nu_x]{|y|}-\E[y\sim \nu_{x+i}]{|y|} + 1\}_i\Cov(\nu_x)
    \\+&\Cov(\nu_x)M_x\Cov(\nu_x) + \Cov(\nu_x)M_x^\top \Cov(\nu_x)\Bigg)\alpha .
\end{align*}
We have the following proposition, which shows the equivalent version of second-order spectral stability.
\begin{proposition}[Second-order spectral stability for linear function]
If $f$ is some pure linear function, then $\text{Second derivative}\leq C(\eta)\cdot \text{First derivative}$ is equivalent to
\begin{equation}\label{eq:SleqF}
   \Pi_x^{1/2}(M_x + M_x^\top)\Pi_x^{1/2} \preceq (1-\eta)C(\eta)I + \diag\{\E[y\sim \nu_{x+i}]{|y|}-\E[y\sim \nu_x]{|y|} - 1\}_i .
\end{equation}
where $\Pi_x = \diag\{\*m(\nu_x)\}^{1/2}$.
\end{proposition}
\begin{example}
    Consider $\Omega = 2^{[n]}$ and the underlying distribution $\mu$ satisfies that $\mu(S) \propto \lambda^{|S|}$. Suppose the pinning $x = \emptyset$ and let $\lambda' = (1-\eta)\lambda$.
    We have that 
    $$
        \*m(\nu_x) = \frac{\lambda'}{1+\lambda'}, \Cov(\nu_x) = \frac{\lambda'}{(1+\lambda')^2} I, \*1^\top (I-\Cov(\nu_{x+i})\Cov(\nu_x)^{-1}) = e_i.
    $$
    $$
        \E[y\sim \nu_{x+i}]{|y|}-\E[y\sim \nu_x]{|y|} - 1 = - \frac{\lambda'}{1+\lambda'}.
    $$
    In this case, Equation \ref{eq:SleqF} is a follows,
    $$
         2I \preceq (1-\eta)C(\eta)\frac{1+\lambda'}{\lambda'}I - I \implies C(\eta)\geq \frac{3\lambda}{1+(1-\eta)\lambda} .
    $$
\end{example}
\end{document}